{\footnotesize
\begin{threeparttable}
\begin{tabular}{lrrrrr}
\toprule
 & & \multicolumn{2}{c}{Within 6 months (\%)} & \multicolumn{2}{c}{Within 12 months (\%)} \\
\cmidrule(lr){3-4}\cmidrule(lr){5-6}
Relative decline & Dismissed & G1 & G3 & G1 & G3 \\
\midrule
All & 24,907,824 & 31.3 & 29.8 & 51.7 & 52.0 \\
5 or more & 5,710,183 & 30.4 & 27.6 & 50.5 & 49.3 \\
10 or more & 2,490,890 & 30.7 & 26.9 & 50.8 & 48.5 \\
15 or more & 1,430,694 & 31.1 & 26.4 & 51.1 & 48.2 \\
20 or more & 839,879 & 30.7 & 25.1 & 50.6 & 46.2 \\
\bottomrule
\end{tabular}
\begin{tablenotes}[flushleft]\footnotesize
\item \textit{Note:} Workers dismissed in regions whose rest grew in the year of separation, six cohorts pooled. Each row pools workers whose origin sector grew at least the indicated number of percentage points less than the rest of the intermediate region (\hyperref[eq:choque_leave_one_out]{Equation~(\ref*{eq:choque_leave_one_out})}); the ``All'' row imposes no restriction. G1 and G3 are the high- and low-portability groups. ``Dismissed'' adds up the three groups. Source: RAIS, authors' calculations.
\end{tablenotes}
\end{threeparttable}
}
